% =====================================================================
% config/packages.tex
% PT: Pacotes e formatação global (fontes, margens, espaçamento, títulos,
%     legendas, bibliografia e acrónimos). Normalmente não precisa de
%     alterar este ficheiro.
% EN: Packages and global formatting (fonts, margins, spacing, headings,
%     captions, bibliography and acronyms). You normally do not need to
%     change this file.
% =====================================================================

% ---------------------------------------------------------------------
% Codificação e idioma / Encoding and language
% ---------------------------------------------------------------------
\usepackage[T1]{fontenc}
\usepackage[utf8]{inputenc}
% PT: A última língua da lista é a língua principal do documento.
%     Para escrever a dissertação em português, troque para
%     \usepackage[english,portuguese]{babel}.
% EN: The last language in the list is the main language of the document.
%     To write the dissertation in Portuguese, change it to
%     \usepackage[english,portuguese]{babel}.
\usepackage[portuguese,english]{babel}

% ---------------------------------------------------------------------
% Fontes / Fonts
% ---------------------------------------------------------------------
% PT: O modelo ESTSetúbal/UNIPS usa Times New Roman no texto e Arial nos
%     títulos. O newtx dá a Times; se não estiver instalado, usa-se o
%     mathptmx. O helvet dá a Helvetica, com as mesmas métricas da Arial.
% EN: The ESTSetúbal/UNIPS model uses Times New Roman for body text and
%     Arial for headings. newtx provides Times; if it is not installed,
%     mathptmx is used instead. helvet provides Helvetica, which has the
%     same metrics as Arial.
% PT: O amsmath e o amsthm têm de ser carregados antes do newtxmath, e o
%     amssymb só com o mathptmx (o newtxmath já tem esses símbolos e dá
%     erro "Command \Bbbk already defined" se os dois forem carregados).
% EN: amsmath and amsthm must be loaded before newtxmath, and amssymb only
%     with mathptmx (newtxmath already has those symbols and fails with
%     "Command \Bbbk already defined" if both are loaded).
\usepackage{amsmath,amsthm}
\IfFileExists{newtxtext.sty}
  {\usepackage{newtxtext,newtxmath}}
  {\usepackage{mathptmx}\usepackage{amssymb}}
\usepackage{helvet}

% ---------------------------------------------------------------------
% Página e espaçamento / Page layout and spacing
% ---------------------------------------------------------------------
% PT: A4 com margens de 3 cm e entrelinha de 1,5.
% EN: A4 with 3 cm margins and 1.5 line spacing.
\usepackage[a4paper,margin=3cm]{geometry}
\usepackage{setspace}
\usepackage{microtype}
\usepackage{csquotes}
% PT: Indenta também o primeiro parágrafo de cada secção.
% EN: Also indent the first paragraph of each section.
\usepackage{indentfirst}

% ---------------------------------------------------------------------
% Figuras, tabelas e listas / Figures, tables and lists
% ---------------------------------------------------------------------
\usepackage{graphicx}
\usepackage{float}
\usepackage{caption}
\usepackage{subcaption}
\usepackage{booktabs}
\usepackage{tabularx}
\usepackage{longtable}
\usepackage{array}
\usepackage{multirow}
\usepackage{siunitx}
\usepackage{enumitem}
\usepackage{xcolor}
\usepackage{pdfpages}
\usepackage{etoolbox}
% PT: TikZ desenha diagramas diretamente em LaTeX (ver capítulo 3).
% EN: TikZ draws diagrams directly in LaTeX (see chapter 3).
\usepackage{tikz}
\usetikzlibrary{arrows.meta,positioning}

% ---------------------------------------------------------------------
% Matemática / Mathematics
% ---------------------------------------------------------------------
% PT: amsmath, amsthm e símbolos AMS já foram carregados com as fontes.
% EN: amsmath, amsthm and AMS symbols were already loaded with the fonts.
\usepackage{mathtools}

% ---------------------------------------------------------------------
% URLs, ligações e referências cruzadas / URLs, links and cross-references
% ---------------------------------------------------------------------
% PT: O xurl não é carregado de propósito: permite quebras após qualquer
%     carácter e parte os DOI a meio de um número. As quebras de URL na
%     bibliografia são afinadas mais abaixo.
% EN: xurl is deliberately not loaded: it allows a break after any
%     character and splits DOIs in the middle of a number. URL breaks in
%     the bibliography are tuned further below.
\PassOptionsToPackage{hyphens}{url}
\usepackage{url}
\usepackage[hidelinks]{hyperref}
\usepackage{bookmark}
% PT: \cref{...} escreve automaticamente "Figure 1.1", "Table 2.3", etc.
%     (com maiúscula, por causa da opção capitalise).
%     O cleveref não tem nomes em português europeu; em texto português
%     escreva o nome à mão: Figura~\ref{...}.
% EN: \cref{...} automatically writes "Figure 1.1", "Table 2.3", etc.
%     (capitalised, because of the capitalise option).
%     cleveref has no European Portuguese names; in Portuguese text write
%     the name by hand: Figura~\ref{...}.
\usepackage[nameinlink,noabbrev,capitalise]{cleveref}

% ---------------------------------------------------------------------
% Acrónimos / Acronyms
% ---------------------------------------------------------------------
% PT: Os acrónimos são definidos em frontmatter/acronyms.tex. A lista só
%     aparece depois de correr o makeglossaries (o latexmk fá-lo sozinho).
% EN: Acronyms are defined in frontmatter/acronyms.tex. The list only
%     appears after running makeglossaries (latexmk does it for you).
% PT: "nomain": só há lista de acrónimos, sem glossário de termos. Para
%     ter também um glossário, retire "nomain", defina os termos com
%     \newglossaryentry e junte \printglossary em main.tex.
% EN: "nomain": there is only an acronym list, no glossary of terms. To
%     also have a glossary, remove "nomain", define the terms with
%     \newglossaryentry and add \printglossary to main.tex.
\usepackage[acronym,nomain,toc]{glossaries}
\makeglossaries

% ---------------------------------------------------------------------
% Bibliografia / Bibliography
% ---------------------------------------------------------------------
% PT: Estilo IEEE (numérico, pela ordem de citação) processado pelo biber.
%     ATENÇÃO: o estilo "ieee" vem do pacote biblatex-ieee, que pode não
%     estar instalado (erro "File 'ieee.bbx' not found"). Instale-o com
%         tlmgr install biblatex-ieee
%     (MiKTeX: instale "biblatex-ieee" no MiKTeX Console) ou use a linha
%     alternativa comentada, que só precisa do biblatex.
% EN: IEEE style (numeric, in order of citation) processed by biber.
%     WARNING: the "ieee" style comes from the biblatex-ieee package, which
%     may not be installed (error "File 'ieee.bbx' not found"). Install it
%     with
%         tlmgr install biblatex-ieee
%     (MiKTeX: install "biblatex-ieee" in the MiKTeX Console) or use the
%     commented alternative line, which only needs biblatex.
\usepackage[backend=biber,style=ieee,sorting=none]{biblatex}
% \usepackage[backend=biber,style=numeric-comp,sorting=none]{biblatex}

% PT: Quebras de URL e DOI na bibliografia: só parte um número em último caso.
% EN: URL and DOI line breaks in the bibliography: break inside a number
%     only as a last resort.
\setcounter{biburllcpenalty}{7000}
\setcounter{biburlucpenalty}{8000}
\setcounter{biburlnumpenalty}{9000}
\appto{\biburlsetup}{\Urlmuskip=0mu plus 1mu\relax}
\AtBeginBibliography{\emergencystretch=4em\relax}

% ---------------------------------------------------------------------
% Espaçamento e numeração / Spacing and numbering
% ---------------------------------------------------------------------
\onehalfspacing
\setlength{\parskip}{0.4em}
\setlength{\parindent}{1.2em}

% PT: O índice mostra até às subsecções; a numeração vai até às subsubsecções.
% EN: The table of contents goes down to subsections; numbering goes down
%     to subsubsections.
\setcounter{tocdepth}{2}
\setcounter{secnumdepth}{3}

\captionsetup{font=small,labelfont=bf}

% PT: Números em formato inglês (1,234.5). Em texto português use
%     \sisetup{output-decimal-marker={,}, group-separator={\,}}
%     (ver o exemplo no capítulo 4).
% EN: Numbers in English format (1,234.5). In Portuguese text use
%     \sisetup{output-decimal-marker={,}, group-separator={\,}}
%     (see the example in chapter 4).
\sisetup{
  group-separator = {,},
  output-decimal-marker = {.}
}

% ---------------------------------------------------------------------
% Títulos (Arial/Helvetica, negrito) / Headings (Arial/Helvetica, bold)
% ---------------------------------------------------------------------
% PT: Títulos em \sffamily (Helvetica = Arial), alinhados à esquerda sem
%     hifenização, como no modelo ESTSetúbal/UNIPS.
% EN: Headings in \sffamily (Helvetica = Arial), ragged right without
%     hyphenation, as in the ESTSetúbal/UNIPS model.
\usepackage{titlesec}
\titleformat{\chapter}[display]
  {\normalfont\sffamily\huge\bfseries}
  {\chaptertitlename\ \thechapter}
  {20pt}
  {\Huge}
\titlespacing*{\chapter}{0pt}{0pt}{40pt}
\titleformat*{\section}{\normalfont\raggedright\sffamily\Large\bfseries}
\titleformat*{\subsection}{\normalfont\raggedright\sffamily\large\bfseries}
\titleformat*{\subsubsection}{\normalfont\raggedright\sffamily\normalsize\bfseries}
\titleformat*{\paragraph}{\normalfont\raggedright\sffamily\normalsize\bfseries}

% ---------------------------------------------------------------------
% Nomes das listas / Names of the lists
% ---------------------------------------------------------------------
% PT: Nomes usados quando a língua principal é o inglês.
% EN: Names used when the main language is English.
\addto\captionsenglish{%
  \renewcommand{\contentsname}{Table of Contents}%
  \renewcommand{\listfigurename}{List of Figures}%
  \renewcommand{\listtablename}{List of Tables}%
  \renewcommand{\bibname}{References}%
  \renewcommand{\refname}{References}%
}
% PT: Se a língua principal for o português, descomente estas linhas.
% EN: If the main language is Portuguese, uncomment these lines.
% \addto\captionsportuguese{%
%   \renewcommand{\contentsname}{Índice}%
%   \renewcommand{\listfigurename}{Índice de Figuras}%
%   \renewcommand{\listtablename}{Índice de Tabelas}%
%   \renewcommand{\bibname}{Referências Bibliográficas}%
%   \renewcommand{\refname}{Referências Bibliográficas}%
% }
